% GENERATED FILE. Edit canonical lessons, then run npm run generate:books.
% canonical-source-hash: fnv1a64:d6cf2f0c411f798a
% canonical-lessons: GE-C427-klimaanlage, GE-C427-heizung, GE-C427-heizkoerper, GE-C427-aufdrehen, GE-C427-zudrehen

\chapter{Warm oder kalt}
\label{ch:ge-warm-oder-kalt}
\hlchaptermodality{427}

\begin{chapteropening}
\textbf{By the end of this chapter:} \emph{I can deal with the temperature in a room: the air conditioning, the heating, the radiator, turning something up and turning it off.}

Complete the last lesson of chapter 427: I can deal with the temperature in a room: the air conditioning, the heating, the radiator, turning something up and turning it off.
\end{chapteropening}

\section[die Klimaanlage]{die Klimaanlage — the air conditioning}

\label{lesson:GE-C427-klimaanlage}



\begin{quote}
Before the new one: say the German for the stay, then the German for to check out.
\end{quote}

\subsection*{You'll want to know: die Klimaanlage}
\textbf{die Klimaanlage} — \textquotedblleft{}the air conditioning\textquotedblright{}.

\begin{quote}
Die Klimaanlage im Zimmer macht viel Lärm.
\end{quote}

\begin{quote}
\emph{The air conditioning in the room makes a lot of noise.}
\end{quote}

Plural: \emph{die Klimaanlagen}. Short: \emph{die Klima}.

\begin{cousinweb}
\emph{Klimaanlage} is built from \emph{Klima} 'climate', from Latin \emph{clima}, from Greek \emph{klíma} 'slope, region' (EtymDB), and \emph{Anlage} 'installation', from \emph{anlegen} 'to lay out' (EtymDB).
\end{cousinweb}

\begin{grammarlens}[title={What you've built}]
One more word for telling, explaining and describing at B1.
\end{grammarlens}

\subsection*{Your turn}
\begin{itemize}
\raggedright
  \item \emph{Say it:} \emph{die Klimaanlage}
  \item \emph{Say it:} \emph{die Klimaanlage}, once more
  \item \emph{Say it:} a sentence of your own with \emph{die Klimaanlage}
  \item \emph{Recall:} say the German for the stay, then the German for to check out, then say \emph{die Klimaanlage} again
\end{itemize}

\begin{tcolorbox}[breakable,colback=teal!4,colframe=teal!35!black,title={Before you move on}]
What does die Klimaanlage mean? (\textquotedblleft{}The air conditioning\textquotedblright{}.) Say it once more.
\end{tcolorbox}
\section[die Heizung]{die Heizung — the heating; the radiator}

\label{lesson:GE-C427-heizung}



\begin{quote}
Before the new one: say the German for to check out, then the German for the air conditioning.
\end{quote}

\subsection*{You'll want to know: die Heizung}
\textbf{die Heizung} — \textquotedblleft{}the heating; the radiator\textquotedblright{}.

\begin{quote}
Im Winter ist die Heizung im Hotel zu schwach.
\end{quote}

\begin{quote}
\emph{In winter the heating in the hotel is too weak.}
\end{quote}

Plural: \emph{die Heizungen}. People often say \emph{die Heizung} for one radiator too.

\begin{cousinweb}
\emph{Heizung} is built on the verb \emph{heizen} 'to heat', inherited from Old High German \emph{heizen} and Proto-Germanic \emph{haitijaną} (EtymDB), a cousin of German \emph{heiß} 'hot'.
\end{cousinweb}

\begin{grammarlens}[title={What you've built}]
One more word for telling, explaining and describing at B1.
\end{grammarlens}

\subsection*{Your turn}
\begin{itemize}
\raggedright
  \item \emph{Say it:} \emph{die Heizung}
  \item \emph{Say it:} \emph{die Heizung}, once more
  \item \emph{Say it:} a sentence of your own with \emph{die Heizung}
  \item \emph{Recall:} say the German for to check out, then the German for the air conditioning, then say \emph{die Heizung} again
\end{itemize}

\begin{tcolorbox}[breakable,colback=teal!4,colframe=teal!35!black,title={Before you move on}]
What does die Heizung mean? (\textquotedblleft{}The heating; the radiator\textquotedblright{}.) Say it once more.
\end{tcolorbox}
\section[der Heizkörper]{der Heizkörper — the radiator}

\label{lesson:GE-C427-heizkoerper}



\begin{quote}
Before the new one: say the German for the air conditioning, then the German for the heating.
\end{quote}

\subsection*{You'll want to know: der Heizkörper}
\textbf{der Heizkörper} — \textquotedblleft{}the radiator\textquotedblright{}.

\begin{quote}
Sie hängt die nassen Socken über den Heizkörper.
\end{quote}

\begin{quote}
\emph{She hangs the wet socks over the radiator.}
\end{quote}

Same form in the plural.

\begin{cousinweb}
\emph{Heizkörper} is built from \emph{heizen} 'to heat', an inherited Germanic verb (EtymDB), and \emph{Körper} 'body', which goes back to Latin \emph{corpus} (Lewis \& Short).
\end{cousinweb}

\begin{grammarlens}[title={What you've built}]
One more word for telling, explaining and describing at B1.
\end{grammarlens}

\subsection*{Your turn}
\begin{itemize}
\raggedright
  \item \emph{Say it:} \emph{der Heizkörper}
  \item \emph{Say it:} \emph{der Heizkörper}, once more
  \item \emph{Say it:} a sentence of your own with \emph{der Heizkörper}
  \item \emph{Recall:} say the German for the air conditioning, then the German for the heating, then say \emph{der Heizkörper} again
\end{itemize}

\begin{tcolorbox}[breakable,colback=teal!4,colframe=teal!35!black,title={Before you move on}]
What does der Heizkörper mean? (\textquotedblleft{}The radiator\textquotedblright{}.) Say it once more.
\end{tcolorbox}
\section[aufdrehen]{aufdrehen — to turn up, to turn on (a tap, the heating)}

\label{lesson:GE-C427-aufdrehen}



\begin{quote}
Before the new one: say the German for the heating, then the German for the radiator.
\end{quote}

\subsection*{You'll want to know: aufdrehen}
\textbf{aufdrehen} — \textquotedblleft{}to turn up, to turn on (a tap, the heating)\textquotedblright{}.

\begin{quote}
Ihm ist kalt, deshalb dreht er die Heizung auf.
\end{quote}

\begin{quote}
\emph{He is cold, so he turns the heating up.}
\end{quote}

Separable: \emph{ich drehe auf}, \emph{ich habe aufgedreht}. \emph{Das Radio aufdrehen} — to turn the radio up.

\begin{cousinweb}
\emph{aufdrehen} is built from \emph{auf} 'up, open' and \emph{drehen} 'to turn', an inherited Germanic verb (EtymDB lists Dutch \emph{draaien} as a cousin).
\end{cousinweb}

\begin{grammarlens}[title={What you've built}]
One more word for telling, explaining and describing at B1.
\end{grammarlens}

\subsection*{Your turn}
\begin{itemize}
\raggedright
  \item \emph{Say it:} \emph{aufdrehen}
  \item \emph{Say it:} \emph{aufdrehen}, once more
  \item \emph{Say it:} a sentence of your own with \emph{aufdrehen}
  \item \emph{Recall:} say the German for the heating, then the German for the radiator, then say \emph{aufdrehen} again
\end{itemize}

\begin{tcolorbox}[breakable,colback=teal!4,colframe=teal!35!black,title={Before you move on}]
What does aufdrehen mean? (\textquotedblleft{}To turn up, to turn on (a tap, the heating)\textquotedblright{}.) Say it once more.
\end{tcolorbox}
\section[zudrehen]{zudrehen — to shut off (a tap, the heating), to turn down}

\label{lesson:GE-C427-zudrehen}



\begin{quote}
Before the new one: say the German for the radiator, then the German for to turn up.
\end{quote}

\subsection*{You'll want to know: zudrehen}
\textbf{zudrehen} — \textquotedblleft{}to shut off (a tap, the heating), to turn down\textquotedblright{}.

\begin{quote}
Vor dem Schlafen dreht sie die Heizung zu.
\end{quote}

\begin{quote}
\emph{Before going to sleep she turns off the heating.}
\end{quote}

Separable: \emph{ich drehe zu}. Its opposite is \emph{aufdrehen}.

\begin{cousinweb}
\emph{zudrehen} is built from \emph{zu} 'shut, closed' and \emph{drehen} 'to turn', an inherited Germanic verb (EtymDB).
\end{cousinweb}

\begin{grammarlens}[title={What you've built}]
One more word for telling, explaining and describing at B1.
\end{grammarlens}

\subsection*{Your turn}
\begin{itemize}
\raggedright
  \item \emph{Say it:} \emph{zudrehen}
  \item \emph{Say it:} \emph{zudrehen}, once more
  \item \emph{Say it:} a sentence of your own with \emph{zudrehen}
  \item \emph{Recall:} say the German for the radiator, then the German for to turn up, then say \emph{zudrehen} again
\end{itemize}

\begin{tcolorbox}[breakable,colback=teal!4,colframe=teal!35!black,title={Before you move on}]
What does zudrehen mean? (\textquotedblleft{}To shut off (a tap, the heating), to turn down\textquotedblright{}.) Say it once more.
\end{tcolorbox}
